\section{Mixed exponent covers from an actual first-profile coefficient}
\label{mx-section}

The two norm exponents in this section may be relatively prime, and
either residual may be a nonunit. The construction controls a finite
list of covers and their coefficient heights. It gives no upper bound
on the heights of their points.

Let $\zeta^2-\zeta+1=0$, $K_0=\mathbb Q(\zeta)$, and
$\gamma=1+\zeta$. Retain the degree-eight splitting field $K$ of
Lemma~\ref{qc-fixed-field}, which contains $K_0$ and the four roots
$\alpha_1,\ldots,\alpha_4$ of
\[
 f(X)=X^4+3X^3+5X^2+3X+1.
\]
Put $\beta_+=-\zeta$ and $\beta_-=-\bar\zeta$, the roots of
$N(X)=X^2+X+1$. These six branch values are distinct: both
polynomials are separable, and $f\bmod N=X^2$, so
$\operatorname{Res}(N,f)=1$.

For positive coprime integers $a,b$, write
$c=a+b$, $x=a/b$, $M=a^2+ab+b^2$, and $F=b^4f(x)$.
Consider genuine oriented profiles
\begin{equation}\label{mx-first-profile}
 a+b\zeta=u\gamma^e v w^h,\qquad
 M=3^eV_0R^h,\qquad V_0=\operatorname N v,\quad R=\operatorname N w\ge7,
\end{equation}
where $u$ is an Eisenstein unit, $e\in\{0,1\}$, $h\ge2$, and
$v,w$ are primitive oriented unramified Eisenstein integers as in
\eqref{lr-profile}. Their oriented supports may overlap. Also suppose
\begin{equation}\label{mx-second-profile}
 F=V_1Q^g,\qquad g\ge9,\quad V_1,Q\in\mathbb Z_{>0},\quad
 Q>1,\quad 1\le V_1<2^g.
\end{equation}
All exponents are integers. The last condition ensures that $V_1$ is
$g$-free. No condition relates $h$ and $g$.

The domain $Q\ge7$ used in Section~\ref{qc-section} follows from
these hypotheses. Indeed, $F$ is odd and $F\equiv1\pmod3$.
The values of $f(r)$ for $r=0,1,2,3,4$ modulo five are
$1,3,2,2,1$. If $5\mid b$, primitivity gives $F\equiv a^4\ne0$
instead. Thus $2,3,5$ do not divide $F$, and $Q>1$ implies $Q\ge7$.

\begin{lemma}[An actual first-profile coefficient]\label{mx-actual-coefficient}
Define
\[
 \xi_0=(u^2\zeta^e v/\bar v)^{-1}.
\]
Then
\begin{equation}\label{mx-first-kummer}
 (w/\bar w)^h=\xi_0\frac{x-\beta_+}{x-\beta_-},
 \qquad h_K(\xi_0)=\frac12\log V_0,
\end{equation}
where $h_K$ denotes absolute logarithmic Weil height. Over all actual
residuals with $V_0\le X$, $X\ge1$, there are at most $CX$
possible $\xi_0$, for an absolute constant $C$. If $v$ is a unit,
the coefficients lie in a fixed finite group of roots of unity,
including when $h$ is divisible by two or three.
\end{lemma}
\begin{proof}
Conjugating \eqref{mx-first-profile}, using $u\bar u=1$ and
$\gamma/\bar\gamma=\zeta$, gives
\[
 \frac{a+b\zeta}{a+b\bar\zeta}
   =u^2\zeta^e\frac{v}{\bar v}
      \left(\frac w{\bar w}\right)^h.
\]
Divide by $b$ and rearrange. The exact oriented-height calculation
of Lemma~\ref{lr-heights} gives
$h_K(v/\bar v)=\tfrac12\log\operatorname N v$:
all archimedean absolute values are one, and the numerator and
denominator ideals are disjoint conjugates. Overlap between the
supports of $v$ and $w$ does not create a conjugate pair within $v$.
Height is unchanged by extending the number field, multiplication
by a root of unity, or inversion. No $h$-th root of $u$ is taken.

Write $v=r+s\zeta$. Since
$\operatorname N v=r^2+rs+s^2\ge3r^2/4,3s^2/4$, the number of
integer coordinate pairs with norm at most $X$ is at most
\[
 \left(\frac{4\sqrt X}{\sqrt3}+1\right)^2\le16X.
\]
The finitely many choices of $u,e$ change only the constant.
This already counts the entire range $V_0\le X$; it is not a
count to be summed again over possible norm values.
\end{proof}

Theorems~\ref{qc-cover-count} and~\ref{qc-coefficient-height}
provide a finite list of $\xi_1,\xi_2,\xi_3\in K^\times$ and
actual coordinates $Z_i\in K^\times$ satisfying
\begin{equation}\label{mx-second-kummer}
 Z_i^g=\xi_i\frac{x-\alpha_i}{x-\alpha_4},\qquad i=1,2,3.
\end{equation}
Here we replace the coefficients of \eqref{qc-covers} by their
inverses; the list sizes and heights are unchanged.

\begin{theorem}[The mixed cover at arbitrary exponents]\label{mx-cover}
Combine \eqref{mx-second-kummer} with
\begin{equation}\label{mx-cover-first-equation}
 Y^h=\xi_0\frac{x-\beta_+}{x-\beta_-}.
\end{equation}
The smooth projective normalization $C_{h,g,\xi}$ is geometrically
connected. Its degree over $\mathbb P^1_x$ and its genus are
\begin{equation}\label{mx-degree-genus}
 D=hg^3,\qquad
 G=1+D\left(2-\frac1h-\frac2g\right)
   =1+2hg^3-g^3-2hg^2.
\end{equation}
Every actual profile \eqref{mx-first-profile}--\eqref{mx-second-profile}
gives a $K$-point of one such curve above $x=a/b$.
\end{theorem}
\begin{proof}
Work over $\bar K$, which contains all required roots of unity
and roots of the nonzero twist coefficients. The three quartic
ratios are independent modulo $g$-th powers even when $g$ is
composite: valuation at $\alpha_i$ is one for its own ratio and
zero for the other two. Their Kummer cover is Galois of degree
$g^3$. Its four branch points are the $\alpha_i$, each of inertia
$g$. At $\alpha_4$ the three simple poles adjoin the same root of
a local uniformizer, so the inertia is diagonal cyclic of order
$g$. Infinity is unramified, as in the proof of
Theorem~\ref{qc-cover-count}.

The first cover is cyclic of degree $h$, by its simple zero at
$\beta_+$, and is branched exactly at $\beta_+,\beta_-$,
each of inertia $h$. Its branch set is disjoint from that of the
quartic cover. The intersection of their function fields over
$\bar K(x)$ is therefore unramified at every point of the
projective line: an intermediate extension can ramify only where
its containing extension ramifies. Riemann--Hurwitz excludes a
nontrivial connected unramified cover of $\mathbb P^1$ in
characteristic zero. Thus this intersection is trivial. Both
extensions are Galois, so their compositum is a field of degree
$hg^3$. This proves geometric connectedness without any
coprimality condition on the exponents.

In the compositum the six ramification contributions sum to
$2D(1-1/h)+4D(1-1/g)$. Infinity remains unramified, and hence
\[
 2G-2=-2D+2D(1-1/h)+4D(1-1/g),
\]
which gives \eqref{mx-degree-genus}. The actual positive rational
$x$ lies outside the six branch values. All four Kummer coordinates
are nonzero, so the affine cover is smooth there in characteristic
zero. The actual point therefore lifts to the normalization over
$K$. This does not assert that every $K$-point arises from a
primitive positive integer seed.
\end{proof}

\begin{theorem}[A mixed residual budget]\label{mx-count-height}
For fixed $h,g,V_1$ as above and all first residuals $V_0\le X$,
$X\ge1$, the actual seeds lift to a list of at most
\begin{equation}\label{mx-fixed-count}
 C_KXg^9 4^{8\omega(V_1)}
\end{equation}
mixed covers. For $L_0\ge0$ and $0\le L_1<\log2$, the union over
$V_0\le\exp(L_0h)$ and $V_1\le\exp(L_1g)$ is covered by at most
\begin{equation}\label{mx-union-count}
 C_Kg^9\exp(L_0h+17L_1g)
\end{equation}
models. The maximum absolute Weil height of an individual
coefficient of the defining equations can be bounded by
\begin{equation}\label{mx-coefficient-bound}
 C_K\bigl(g+\log V_0+\log V_1+1\bigr).
\end{equation}
The same bound, with a different constant, holds for their
projective coefficient height. The degrees in $Y$ and $Z_i$
remain $h$ and $g$, respectively.
\end{theorem}
\begin{proof}
The QC count gives $C_Kg^94^{8\omega(V_1)}$ second-norm
models for fixed $g,V_1$, with
$h_K(\xi_i)\le2\log V_1+C_Kg$. Its proof includes the fixed
field's unit rank three, class-group choices and the actual
depth-one contribution above thirteen. Lemma~\ref{mx-actual-coefficient}
adds at most $CX$ first coefficients for the entire residual range.
There is no additional independent $h^3$ unit-class factor:
the actual first profile has already supplied its precise
coefficient. This proves \eqref{mx-fixed-count}.

Since $V_1\ge2^{\omega(V_1)}$, one has
$4^{8\omega(V_1)}\le V_1^{16}$. Summing over the integers
$V_1\le Z$, $Z\ge1$, costs at most $Z^{17}$. Substituting
$X=\exp(L_0h)$ and $Z=\exp(L_1g)$ proves
\eqref{mx-union-count}. No second summation over first residual
norms is required.

Clear the linear denominators to obtain
\[
 (x-\beta_-)Y^h-\xi_0(x-\beta_+)=0,
 \qquad
 (x-\alpha_4)Z_i^g-\xi_i(x-\alpha_i)=0.
\]
The branch values are fixed, and product heights are bounded by
the sums of factor heights. Equation~\eqref{mx-first-kummer}
and the QC coefficient bound prove
\eqref{mx-coefficient-bound}. There are a fixed number of
nonzero coefficients, so their projective coefficient height
has the same bound up to a constant. This reasoning supplies
no point-height estimate.
\end{proof}

\begin{corollary}[Height separation when the first exponent dominates]
\label{mx-height-separation}
Consider any sequence of actual profiles as above with
\begin{equation}\label{mx-sequence-domain}
 h\longrightarrow\infty,\qquad g/h\longrightarrow0,\qquad
 \frac{\log V_0}{h}\longrightarrow0,\qquad
 \frac{\log V_1}{g}\ \hbox{bounded}.
\end{equation}
For the chosen models, let $H_{\rm coef}$ be the coefficient-height
bound in \eqref{mx-coefficient-bound}, and put
$H_{\rm point}=h_K(a/b)=\log\max(a,b)$. Then
\begin{equation}\label{mx-height-ratio}
 \frac{H_{\rm coef}+\log h+\log g}{H_{\rm point}}
 \longrightarrow0.
\end{equation}
\end{corollary}
\begin{proof}
The hypotheses give $H_{\rm coef}=o(h)$. In the other direction,
\eqref{mx-first-profile} implies $M\ge R^h\ge7^h$, while
$M<c^2$. Consequently
\[
 H_{\rm point}\ge\log c-\log2
   \ge\frac h2\log7-\log2.
\]
Also $\log h=o(h)$ and $\log g=o(h)$ because $9\le g=o(h)$.
Division proves \eqref{mx-height-ratio}.
\end{proof}

This separation allows two nonunit residuals and relatively prime
exponents. It asserts neither the existence of the displayed
sequence nor a contradiction. A sufficient further input would be
an upper bound
\[
 H_{\rm point}\le A\bigl(H_{\rm coef}+\log h+\log g+1\bigr)
\]
for the actual lifted points in this family, with $A$ independent
of the exponents and residuals. Such a point-height theorem is
not proved here. Comparable exponents, larger moving residuals,
and points whose heights exceed the coefficient budget remain
unresolved. The present ordinary arithmetic-geometric proof is
independently reviewed; the Kummer, ramification and height
arguments in this section are not Lean formalizations.
